\section{Incidencia areal--volumétrica y entropía topológica}

El calendario de TPK induce la bisagra entre los modos aditivo y
multiplicativo. En el bloque primitivo
\[
 \tau_3=(+1,-1,0),
\]
la regla
\[
 +1:m\mapsto\Sigma,
 \qquad
 -1:m\mapsto\Pi,
 \qquad
 0:m\mapsto m
\]
produce la órbita cíclica \((\Sigma,\Pi,\Pi)\). El lector geométrico
\[
 \Sigma\longmapsto\mathscr H_{\rm ar},
 \qquad
 \Pi\longmapsto\mathscr H_{\rm vol}
\]
transporta esos modos a las hojas areal y volumétrica.

\begin{theorem}[Incidencia necesaria de las hojas]
\label{thm:incidencia-hojas-anexo}
En el orden \((\mathscr H_{\rm ar},\mathscr H_{\rm vol})\), la matriz
primitiva de incidencias es
\begin{equation}
 \boxed{
 F_{\rm av}=
 \begin{pmatrix}0&1\\1&1\end{pmatrix}.}
 \label{eq:fav-anexo-informacion}
\end{equation}
Los calendarios de longitudes \(9,27,108\) producen
\[
 N_9=3F_{\rm av},
 \qquad N_{27}=9F_{\rm av},
 \qquad N_{108}=36F_{\rm av}.
\]
\end{theorem}

\begin{proof}
Los pares consecutivos del ciclo son
\(\Sigma\to\Pi\), \(\Pi\to\Pi\) y \(\Pi\to\Sigma\). Cada uno aparece una
vez y \(\Sigma\to\Sigma\) no aparece. Esto fija las cuatro entradas de
\(F_{\rm av}\). Los calendarios largos contienen, respectivamente, tres,
nueve y treinta y seis copias del bloque primitivo.
\end{proof}

\begin{theorem}[Ley de Fibonacci y medida de máxima entropía]
\label{thm:fibonacci-parry-anexo}
Para \(n\geq1\),
\begin{equation}
 F_{\rm av}^{n}
 =\begin{pmatrix}F_{n-1}&F_n\\F_n&F_{n+1}\end{pmatrix}.
 \label{eq:potencias-fav-anexo}
\end{equation}
El subshift de memoria uno definido por \(F_{\rm av}\) tiene
\begin{equation}
 \boxed{h_{\rm top}=\log\varphi},
 \qquad
 \boxed{\frac{\mu_{\rm ar}}{\mu_{\rm vol}}=\varphi^{-2}}.
 \label{eq:entropia-parry-hojas-anexo}
\end{equation}
\end{theorem}

\begin{proof}
La identidad matricial se obtiene por inducción a partir de
\(F_{n+1}=F_n+F_{n-1}\). La raíz de Perron de \(F_{\rm av}\) es \(\varphi\),
de donde \(h_{\rm top}=\log\varphi\). Como la matriz es simétrica, los pesos
de Parry \cite{Parry1964} son proporcionales al cuadrado de un vector de Perron
\((1,\varphi)\), es decir, a \((1,\varphi^2)\).
\end{proof}

Las expresiones \(\log3\) y \(\log\varphi\) quedan así separadas por su
genealogía. La primera cuantifica una elección local uniforme entre tres
estados; la segunda mide el crecimiento de las historias permitidas por la
incidencia de hojas. La relación areal--volumétrica es una propiedad del
retorno completo y por ello aparece con el exponente dos.

